% Source MD: E:\wanru\PRML_Course\NoteMD\Week1\第1节课 机器学习基本概念.md
\documentclass[main.tex]{subfiles}

\graphicspath{{\subfix{figures/}}}

\begin{document}

\chapter{机器学习基本概念}
\label{chap:w01l01-ml-intro}

\begin{keypointbox}
机器学习的核心任务是从资料中找出一个足够好的函数。
以课程中的 YouTube 点阅人数预测为例，本节重点是三件事：先写出带参数的模型，再定义损失函数，最后用优化算法把参数找出来。
\end{keypointbox}

\section{什么是机器学习}
\label{sec:w01l01-what-is-ml}

\begin{definition}{机器学习}{w01l01-machine-learning}
机器学习研究的是如何利用资料自动找出一个函数，使它能够把输入映射到合理的输出。
\end{definition}

课程中的一句话总结是：机器学习就是让机器具备“找函数”的能力。
当任务足够复杂时，例如语音辨识、影像辨识或围棋落子，人类往往无法直接手写出精确规则，
这时就希望机器能从数据中自动学出这个复杂函数。

\section{机器学习的任务类别}
\label{sec:w01l01-task-types}

根据输出物的型态不同，课程中将机器学习任务粗略分成三类：
\begin{enumerate}
  \item \textbf{Regression（回归）}：输出是一个数值。例如根据今天和过去几天的空气品质指标，预测明天中午的 PM2.5。
  \item \textbf{Classification（分类）}：输出是在有限类别里做选择。例如垃圾邮件过滤，或从棋盘上 $19 \times 19$ 个位置中选择下一手。
  \item \textbf{Structured Learning（结构化学习）}：输出具有结构，例如句子、文档、图像等，更接近让机器“创造”一个复杂物件。
\end{enumerate}

\begin{examplebox}
分类问题像在做选择题，回归问题像在估计一个连续数值，结构化学习则要求模型直接生成一个有组织的结果。
\end{examplebox}

\section{找函数的三个步骤}
\label{sec:w01l01-three-steps}

课程用预测 YouTube 频道隔天总观看次数作为例子，说明机器学习如何一步步找出函数。

\subsection{Step 1：写出带未知参数的模型}
\label{sec:w01l01-model}

第一步是先根据领域知识写出一个有未知参数的函数。最简单的猜测是：隔天的观看人数和前一天的观看人数有关，
于是可以先写成
\begin{equation}
  y = b + w x_1.
  \label{eq:w01l01-linear-model}
\end{equation}

之所以称 \cref{eq:w01l01-linear-model} 为线性模型，是因为输出 $y$ 由各个输入特征的线性组合再加上一个常数项得到。
固定 $b$ 不变时，$x_1$ 每增加 1，预测值就会改变 $w$；固定 $w$ 不变时，$b$ 只会让整条直线整体上移或下移。
换句话说，$w$ 决定了模型对输入变化的敏感程度，$b$ 决定了模型在没有任何特征讯号时的基准预测。

这里各个符号的意义是：
\begin{itemize}
  \item $y$：要预测的目标，也就是隔天的观看人数。
  \item $x_1$：已知资讯，也就是前一天的观看人数，称为特征（feature）。
  \item $w$：权重（weight）。
  \item $b$：偏差（bias）。
\end{itemize}

像 \cref{eq:w01l01-linear-model} 这种带有未知参数的函数，就叫做模型（model）。

\begin{examplebox}
如果某频道昨天的观看人数为 $x_1=100$，并且我们暂时取 $w=1.1$、$b=20$，那么模型会给出预测
\[
  y = 20 + 1.1 \times 100 = 130.
\]
若保持 $w=1.1$ 不变，只把 $b$ 改成 $50$，预测就会整体抬高到 $160$；若保持 $b=20$ 不变，把 $w$ 改成 $0.8$，
则模型对昨天观看人数的反应会变弱，预测会变成 $100$。这个例子对应了 weight 与 bias 的两种不同作用。
\end{examplebox}

\subsection{Step 2：定义损失函数}
\label{sec:w01l01-loss}

第二步要定义损失函数（loss），衡量某一组参数到底好不好。
做法是把训练资料代入模型，得到预测值 $y$，再跟真实标签 $\hat{y}$ 比较。

单笔资料的误差常见写法包括
\begin{align}
  e_{\mathrm{MAE}} &= \abs{y - \hat{y}}, \label{eq:w01l01-mae} \\
  e_{\mathrm{MSE}} &= \paren{y - \hat{y}}^2. \label{eq:w01l01-mse}
\end{align}

\subsubsection{从单点误差到整体 Loss}

\cref{eq:w01l01-mae,eq:w01l01-mse} 先描述的是单一训练样本的误差。
这一步只回答“模型在这一笔资料上错了多少”，还没有回答“这组参数在整份训练资料上整体表现如何”。
因此我们要把每一笔样本的误差收集起来，再做平均，才得到真正用于训练的 loss。

把所有训练样本的误差加总并平均，就得到整体损失：
\begin{equation}
  L = \frac{1}{N}\sum_{n=1}^{N} e_n.
  \label{eq:w01l01-loss}
\end{equation}

\cref{eq:w01l01-loss} 中的 $N$ 是训练样本数，$e_n$ 是第 $n$ 笔资料对应的误差。
取平均而不是单纯加总，有两个好处：第一，不同大小的数据集可以直接比较 loss；第二，loss 的尺度不会随着资料笔数线性放大。

\begin{examplebox}
假设有三笔训练资料，它们的预测误差分别是 $10$、$-5$、$20$。那么：
\begin{align*}
  \text{MAE} &= \frac{\abs{10} + \abs{-5} + \abs{20}}{3}
  = \frac{35}{3}, \\
  \text{MSE} &= \frac{10^2 + (-5)^2 + 20^2}{3}
  = \frac{525}{3} = 175.
\end{align*}
这个 toy example 说明了一件事：MSE 会对大误差给出更重的惩罚，因此当少数样本误差特别大时，MSE 往往会比 MAE 更敏感。
\end{examplebox}

\subsubsection{Loss 的几何意义}

\begin{intuitionbox}
损失函数的值越小，代表这组参数在训练资料上的表现越好。若模型只有一个参数，就可以把 $L(w)$ 画成一条曲线；若模型有两个参数，
就可以把 $L(w,b)$ 画成一张曲面或等高线图。课程里反复提到的 error surface，本质上就是“参数位置”与“loss 高低”之间的对应关系。
\end{intuitionbox}

从这个角度看，训练的目标就很直观：我们在参数空间中移动，希望走到 error surface 上更低的位置。
所谓“找到好模型”，其实就是“找到一组让曲面高度尽量低的参数”。

\subsection{Step 3：做优化}
\label{sec:w01l01-optimization}

第三步是解一个优化问题，找出最好的参数：
\begin{equation}
  w^\star, b^\star = \argmin_{w,b} L.
  \label{eq:w01l01-opt-target}
\end{equation}

符号 $\argmin$ 的意思不是直接给出最小的 loss 值，而是找出“让 loss 最小的那组参数”。
因此 \cref{eq:w01l01-opt-target} 的输出是参数 $w^\star,b^\star$，而不是一个数字。

课程中最核心的方法是梯度下降法（gradient descent）。
如果先只考虑一个参数 $w$，更新规则是
\begin{equation}
  w^{t+1} = w^t - \eta \frac{\partial L}{\partial w}\bigg\vert_{w=w^t},
  \label{eq:w01l01-gd-update}
\end{equation}
其中 $\eta$ 是学习率（learning rate），属于需要人工设定的超参数（hyperparameter）。

\subsubsection{为什么要减去梯度}

\cref{eq:w01l01-gd-update} 中最关键的符号是负号。它来自一阶导数的方向意义：
\begin{itemize}
  \item 若 $\frac{\partial L}{\partial w} > 0$，表示当 $w$ 往右增加时，loss 会先上升，因此我们应该把 $w$ 往左调小。
  \item 若 $\frac{\partial L}{\partial w} < 0$，表示当 $w$ 往右增加时，loss 会先下降，因此我们应该把 $w$ 往右调大。
\end{itemize}

这两种情况都可以统一写成“沿着负梯度方向移动”，也就是从当前点朝着最陡下降的方向迈一步。

\begin{algorithm}[htbp]
  \caption{单参数梯度下降}
  \label{alg:w01l01-gradient-descent}
  \KwInput{初始参数 $w^0$，学习率 $\eta$，最大更新次数 $T$}
  \KwOutput{更新后的参数 $w^T$}
  \For{$t \leftarrow 0$ \KwTo $T-1$}{
    计算梯度 $g^t \leftarrow \dfrac{\partial L}{\partial w}\big\vert_{w=w^t}$\;
    更新参数 $w^{t+1} \leftarrow w^t - \eta g^t$\;
  }
  \KwRet{$w^T$}
\end{algorithm}

如果模型有两个参数 $w$ 和 $b$，则分别对两个参数求偏导并更新：
\begin{align}
  w^{t+1} &= w^t - \eta \frac{\partial L}{\partial w}\bigg\vert_{w=w^t,b=b^t}, \label{eq:w01l01-gd-w} \\
  b^{t+1} &= b^t - \eta \frac{\partial L}{\partial b}\bigg\vert_{w=w^t,b=b^t}. \label{eq:w01l01-gd-b}
\end{align}

\subsubsection{学习率与更新行为}

学习率 $\eta$ 决定了每次更新要跨多大一步。它本身不来自数据学习，而是由使用者手动设定。
\begin{itemize}
  \item 若 $\eta$ 太大，参数可能一次跳过最低点，在曲面两侧来回震荡，甚至直接发散。
  \item 若 $\eta$ 太小，loss 虽然会下降，但下降速度很慢，训练会非常耗时。
\end{itemize}

因此，梯度告诉我们“该往哪个方向走”，学习率则决定“这一步走多远”。

\subsubsection{单参数与多参数的关系}

\cref{eq:w01l01-gd-update} 是单参数版本，\cref{eq:w01l01-gd-w,eq:w01l01-gd-b} 则是多参数版本。
本质上两者是同一个想法：只不过当参数变多时，我们不再更新一条数轴上的位置，而是在更高维的参数空间里，同时沿着各个坐标方向做调整。
把所有参数合并成向量 $\theta$ 以后，这个更新规则就会变成熟悉的统一形式
\[
  \theta^{t+1} = \theta^t - \eta \grad L\paren{\theta^t}.
\]

\begin{warningbox}
当梯度变成零时，训练会停下来。这不一定代表已经找到全域最佳解；模型也可能停在局部极小值，甚至是后续课程要讨论的鞍点附近。
\end{warningbox}

\section{线性模型的改进}
\label{sec:w01l01-linear-model}

如果只看前一天的资料，模型通常太简单。进一步利用资料中的周期性，例如一周为一个循环，可以把模型扩展为
\begin{equation}
  y = b + \sum_{j=1}^{7} w_j x_j.
  \label{eq:w01l01-7day-model}
\end{equation}

虽然 \cref{eq:w01l01-7day-model} 看起来比单变量版本复杂得多，但它仍然是线性模型。
原因在于：模型的输出依旧是各个特征乘上权重后做加总，中间没有额外的非线性变换，也没有特征之间的乘积项。
换句话说，它只是从“看前一天”升级成了“同时看前七天”，并没有改变模型的基本型态。

更一般地，凡是把特征乘上权重再加总、最后加上偏差的模型，都属于线性模型（linear model）。
增加观察天数往往会让训练误差下降，但测试误差不一定同步改善，这也为后面的泛化与过拟合问题埋下伏笔。

\begin{summarybox}
本节建立了课程最基础的训练框架：
\begin{enumerate}
  \item 先用领域知识写出带未知参数的模型；
  \item 再定义损失函数衡量参数好坏；
  \item 最后用梯度下降去找最小损失对应的参数。
\end{enumerate}
后续所有更复杂的模型，本质上都还是沿着这三步展开。
\end{summarybox}

\end{document}
